\scanpage{297}
\mspage{38}
\deleted{portée est contenue dans $[f(0_{\unsure{\tau|t}}),f(1_\tau)]$, et est égale
à cet ensemble si $f$ est croissante.}

Mais $J\simeq\unsure{\mathbb{\Delta}|\Delta}_i^\varepsilon$, et on a donc un \deleted{application}\inserted{morphisme}
\uline{surjectif}
\[
 f_0\colon\mathbb{\Delta}_i^\varepsilon\longrightarrow\deleted{\unsure{\mathbb{\Delta}_1}}J'\simeq\mathbb{\Delta}_1^\varepsilon
\]
On voit tout de suite qu'il existe alors une
\uline{unique} flèche de \deleted{\unsure{de $f$}} transition $u$ de
$\mathbb{\Delta}_i^\varepsilon$, telle que $f(u)\ne\mathrm{id}$. Donc on trouve
le corollaire suivant, plus précis que le
cor. 4

\uline{Cor 5} Soit $f\colon I_\tau\unsure{\mathbin{-}|\longrightarrow}I_{\tau'}$ un morphisme
croissant d'intervalles orientés, \deleted{et soit}\inserted{avec} $I_{\tau'}$
réduit, et soit $\mathrm{Tr}(\tau,f)$ l'ens des flèches de
transition \inserted{$\underline{u}_i$} de $I_\tau$ telles que $f(\underline{u}_i)\ne\mathrm{id}$ i.e. telles
que \deleted{la \unsure{portée|pointu}} $\pi(u_i,f)=1$. Alors $f$ induit
une application
\[
\mathrm{Tr}(\tau,f)\longrightarrow\mathrm{Tr}(\tau')
\]
qui est bijective, et (bien sûr) \deleted{croissante}
donc un isom. d'ens. ordonnés.

Ainsi l'application $\mathrm{Tr}(\tau')\unsure{\mathbin{-}|\longrightarrow}\mathrm{Cr}(\tau,f)$ de
tantôt se précise par une application bijective
\[
\mathrm{Tr}(f)^*\colon\mathrm{Tr}(\tau')\longrightarrow\mathrm{Tr}(\tau,f),
\]
reliée à $\mathrm{Cr}(f)^*$ plus haut par le triangle
commutatif
\diagram{\begin{tikzcd}[ampersand replacement=\&,column sep=huge,row sep=small]
\& \mathrm{Tr}(\tau,f)\arrow[dd,Rightarrow,"{\unsure{\sim|\mathord{2}}\ \deleted{\unsure{\varepsilon|E}}\alpha}"]\\
\mathrm{Tr}(\tau')\arrow[ur,"{\mathrm{Tr}(f)^*}"]\arrow[dr,"{\mathrm{Cr}(f)^*}"'] \& {}\\
\& \mathrm{Cr}(\tau,f)
\end{tikzcd}},
% archive sheet number: 296
